\documentclass[10pt]{article}
\usepackage[a4paper,margin=1.8cm]{geometry}
\usepackage{booktabs}
\usepackage{array}
\usepackage{xcolor}
\usepackage{amsmath}
\usepackage{amssymb}
\usepackage{siunitx}
\usepackage{longtable}
\usepackage{graphicx}
\graphicspath{{../}{./}}
\usepackage{hyperref}

% ---- source category tags ------------------------------------------------
\definecolor{cgen}{RGB}{20,80,170}   % genet23 (verified against the PDF)
\definecolor{ccar}{RGB}{0,130,120}   % caruel13 fallback
\definecolor{ccode}{RGB}{120,70,160} % authors' code (Refs 33/6)
\definecolor{csvz}{RGB}{90,90,90}    % svZeroDSolver numerics/plumbing
\definecolor{cnone}{RGB}{190,40,40}  % not found in either paper
\newcommand{\gen}{\textcolor{cgen}{\textbf{genet23}}}
\newcommand{\car}{\textcolor{ccar}{\textbf{caruel13}}}
\newcommand{\code}{\textcolor{ccode}{\textbf{authors-code}}}
\newcommand{\svz}{\textcolor{csvz}{\textbf{svZeroD}}}
\newcommand{\none}{\textcolor{cnone}{\textbf{NOT-in-papers}}}

\newcolumntype{P}[1]{>{\raggedright\arraybackslash}p{#1}}

\title{\textbf{ChamberCylinder reproduction of Genet~(2023):\\ parameter \& input-function provenance (verified against the PDF)}}
\author{svZeroDSolver \texttt{ChamberCylinder} block}
\date{\today}

\begin{document}
\maketitle

\noindent
Every parameter and input function of the \texttt{ChamberCylinder} reproduction of
Genet~et~al.~(2023, \emph{Int.\ J.\ Numer.\ Meth.\ Biomed.\ Engng.}\ e3711), verified
against the paper. Equation and Table numbers refer to that paper; the reduced
active law and circulation trace back to Caruel~et~al.~(2013, \emph{BMMB}
\textbf{13}:897).

\paragraph{Source legend.}
\begin{tabular}{ll}
\gen & value/form verified in genet23 Table~1 or its equations \\
\car & not tabulated in genet23; taken from Caruel~(2013) \\
\code & function from the authors' code \texttt{cardiac\_input\_functions.py} (genet23 Refs 33/6) \\
\svz & svZeroDSolver numerical or plumbing choice (not a physical parameter) \\
\none & appears in genet23 but is \emph{not} faithfully reproduced (see notes) \\
\end{tabular}

% =========================================================================
\section*{Geometry (Table~1) --- all \gen}
\begin{longtable}{l l l P{6.6cm}}
\toprule Parameter & Value & Source & Note \\ \midrule\endhead
$R_i,\,R_e,\,L$ & 19.035, 32.969, \SI{57.104}{mm} & \gen & derived from Table~1's chamber \SI{65}{mL}, wall \SI{130}{mL}, aspect 1.5; Table~1 prints them rounded (19.0, 33.0, 57.1) --- see audit A6 \\
$\alpha_i/\alpha_e$ (Eq.~22) & $+60^\circ/-60^\circ$ & \gen & linear helix angle $\alpha(R)$ across the wall \\
\bottomrule
\end{longtable}

% =========================================================================
\section*{Passive material (Eqs.~25, 26, 28; Table~1) --- all \gen}
\begin{longtable}{l l l P{6.6cm}}
\toprule Parameter & Value & Source & Note \\ \midrule\endhead
$W^e$ form (Eq.~25) & --- & \gen & $C_1(\bar I_1{-}3){+}C_2(\bar I_2{-}3){+}C_3e^{C_4(\bar I_1-3)^2}{+}C_5e^{C_6(\bar I_4-1)^2}$; matches kernel exactly \\
$C_1{,}C_2{,}C_3{,}C_4{,}C_5{,}C_6$ & 7,0,700,2,50,4 & \gen & Table~1 (Pa where dimensional) \\
$\gamma$ (viscosity, Eq.~28) & \SI{70}{Pa\,s} & \gen & $W^v=\tfrac{\gamma}{2}|\dot E|^2\Rightarrow\Sigma^v=\gamma\dot E$ (full Kelvin--Voigt); matches kernel \\
incompressibility (Eqs.~26--27) & mixed u/p & \gen & Lagrange multiplier $p$, $\int(J{-}1)\hat p=0$ \\
$\kappa$ (penalty) & $10^7$ & \svz & fallback only (mixed off); unused \\
\bottomrule
\end{longtable}

% =========================================================================
\section*{Active contraction (Eqs.~29--33; Table~1)}
\begin{longtable}{l l l P{6.6cm}}
\toprule Parameter & Value & Source & Note \\ \midrule\endhead
$\sigma_0$ & \SI{65}{kPa} & \gen & Table~1 \\
$k_0$      & \SI{260}{kPa} & \gen & Table~1 \\
$k_s$      & $10^8$ & \gen & Table~1 \\
$\mu$      & \SI{70}{Pa\,s} & \gen & Table~1 (fiber dashpot, Eq.~32) \\
$\alpha$ (\texttt{bcs\_alpha}) & \textbf{12.0} & \gen & \textbf{Table~1 ``time constant''} (Eq.~33 $\alpha|\dot e_c|$) --- \emph{not} a guess \\
\texttt{active\_stretch\_pow} & 1 & \gen & Eq.~30: $\sigma_{1D}=T_\text{fib}/(1{+}e_\text{fib})$, $1{+}e_\text{fib}=\sqrt{I_4}$ (Eq.~59 as printed would give $p{=}2$; audit B2) \\
$|\dot e_c|$ smoothing & $10^{-10}$ & \svz & $\sqrt{\dot e_c^2+\epsilon}$ in Eq.~33; numerically inert (audit A9) \\
relaxation law (Eq.~33) & fixed ($w{=}1$) & \gen & Eq.~33 decay $=|\nu|+\alpha|\dot e_c|$, \textbf{no length-dependent $w$} $\Rightarrow$ \texttt{bcs\_relax=off} is correct \\
$n_0(e_c)$ & 9-pt PWL & \code & Table~1 cites Ref~6; the authors' \texttt{starling} curve (matches exactly). Large, correct effect (endo $1.0\to$ epi $0.4$) \\
\bottomrule
\end{longtable}

% =========================================================================
\section*{Input functions $\nu(t)$, $P_\text{at}(t)$ (Refs.~33/6, authors' code)}
\begin{longtable}{l l l P{6.6cm}}
\toprule Parameter & Value & Source & Note \\ \midrule\endhead
$\nu(t)$ (\texttt{nagumo}) & --- & \code & ECG-parametric: delay 0.1, QRSd 0.08, QTc 0.33 (Bazett $QT{=}QTc\sqrt{T}$), ramp 0.01; $\alpha_\text{max/min}{=}\pm35$. Kernel port matches to $3\times10^{-14}$ \\
$P_\text{at}(t)$ & 4.5 / 7 mmHg & \code & 603 / \SI{938}{Pa}; kick at $T{-}PQ{+}$ramp, $T{-}AV{+}$ramp (PQ 0.14, AV 0.08) \\
period $T$ & \SI{0.8}{s} & \gen & \\
\bottomrule
\end{longtable}
\noindent\footnotesize The ECG constants live in the authors' code (implementing genet23 Refs 33/6), not in genet23 Table~1. \normalsize

% =========================================================================
\section*{Valves (Eq.~35) \& two-stage Windkessel (Eq.~36; Table~1)}
The valve law (Eq.~35) uses positive parts $\langle\cdot\rangle_+$, i.e.\ a sharp
switch --- so the svZeroD \texttt{PiecewiseValve} diode ($R=1/K$) is \emph{faithful},
not an approximation. Mapping Eq.~36: mitral $R_{\min}{=}K_{at}^{-1},R_{\max}{=}K_{iso}^{-1}$
reproduces the atrial terms $K_{at}\langle P_{at}{-}P_v\rangle_+ - K_{iso}\langle P_v{-}P_{at}\rangle_+$ exactly.
The Windkessel is a chain of \texttt{BloodVessel} blocks,
$C_\text{ar}$ at the aortic-valve outlet, then $R_p$ with $C_d$, then $R_d$ to $P_\text{vs}$,
which is Eq.~36b--c exactly (audit A1; the flow-through \texttt{BloodVesselRC} used earlier was not).
\begin{longtable}{l l l P{6.4cm}}
\toprule Parameter & Value & Source & Note \\ \midrule\endhead
$K_{at}^{-1},K_{ar}^{-1},K_{iso}^{-1}$ & 1.1111e5, 7.6923e4, 2.0e9 & \gen & Table~1 (Pa\,s\,m$^{-3}$) \\
$R_p,\,C_\text{ar}$ & 1.35e7, 2.115e-10 & \gen & Table~1 (proximal Windkessel) \\
$R_d,\,C_d$ & 1.0949e8, 2.0158e-8 & \gen & Table~1 (distal Windkessel) \\
aortic $R_{\max}$ & $10^{12}$ ($\approx\infty$) & \gen & Eq.~36 arterial line has \emph{no} closed-leak; \emph{corrected} from $K_{iso}^{-1}$ (removed spurious regurgitation) \\
$C_\text{valve}$ & 9e-9 & \none & Table~1; Eq.~36 adds $C_\text{valve}\dot P_v$ at the ventricular node. Faithfully applied, \SI{9e-9}{m^3/Pa} ($\approx$\SI{1.2}{mL/mmHg}, $43\times C_\text{ar}$) is a huge compliance: it absorbs $\sim$\SI{90}{mL} over the pressure swing, collapsing the isovolumic phase (ESV 39 vs.\ 74, peak $P$ 8.5 vs.\ 12.8) --- inconsistent with the paper's own Fig.~5. Only $\le$9e-11 is negligible, so we run $C_\text{valve}{\approx}0$ (see Discrepancies) \\
$P_\text{vs}$ (venous) & \SI{350}{Pa} & \textcolor{cnone}{\textbf{FITTED}} & \textbf{not in genet23 Table~1} (only in Eq.~36). \emph{Fitted} to the Fig.~5 ESV: \SI{350}{Pa} ($\approx$\SI{2.6}{mmHg}) gives ESV \SI{74}{mL}; Caruel's \SI{1000}{Pa} overshoots to \SI{76.4}{mL}. Sets only the ESV/afterload floor ($\sim$\SI{1.8}{mL}/\SI{500}{Pa}); EDV/peak-$P$/twist insensitive. The one remaining fitted quantity \\
\bottomrule
\end{longtable}

% =========================================================================
\section*{Numerics}
\begin{longtable}{l l l P{6.4cm}}
\toprule Parameter & Value & Source & Note \\ \midrule\endhead
time step $\Delta t$ & \SI{1}{ms} & \gen & $T/800$, matches genet23 \S3.1 \\
interpolation / quadrature & $n_e\times$P1/P0, 3-pt Gauss & \svz & genet23: 1 element, P3 displacement / P2 pressure, quadrature order 5. See the convergence section \\
time integrator & generalized-$\alpha$, $\rho_\infty{=}0$ & \svz & 2nd order, L-stable; genet23 uses the energy-consistent midpoint scheme (Eqs.~51--65). See the convergence section \\
rigid-body pins & static only & \gen & $\varphi(R_i){=}\eta(R_i){=}0$ only without inertia (\S2.4); free in dynamics (audit A3) \\
cycles & 12, take last & \svz & periodic to 0.01\,mL (cycle 10: 0.06\,mL); genet23 plots cycle 5 of its own start-up \\
initial state & reference, all $P{=}P_\text{at}$ & \svz & genet23 preloads statically to $P_\text{at}$, arterial start unstated; our cycle 5 is still 2.3\,mL (ESV) from periodic \\
\bottomrule
\end{longtable}

% =========================================================================
\clearpage
\section*{Comparison with the digitised Genet~(2023) figures}
\noindent Each figure overlays the ChamberCylinder result (dashed, coloured by variation) on the
digitised Genet trace (solid; stored as CSV in \texttt{scripts/}, extracted from the published
figures --- no copyrighted images in the repo). All runs use the model after the formulation
audit (identical to the paper's continuous equations, with the \emph{full} Table-1 damping
$\gamma{=}\mu{=}\SI{70}{Pa\,s}$), $n_e=6$, $\Delta t=\SI{1}{ms}$ and the last of 12 cycles.
Layout, axes and units follow the paper: cavity pressure
(with the twin kPa/mmHg axis), cavity volume, PV loop, and ventricular torsion. Pressure, volume
and the PV loop match across every case; the torsion untwist is discussed at the end.

\newcommand{\figone}[2]{%
  \begin{center}\textbf{#2}\\[1pt]\includegraphics[width=0.88\linewidth]{#1.png}\end{center}\vspace{4pt}}

\figone{fig5}{Fig.~5 --- baseline beat}
\figone{fig6}{Fig.~6 --- aspect ratio $L/2R_i$}
\clearpage
\figone{fig7}{Fig.~7 --- internal volume}
\figone{fig8}{Fig.~8 --- wall volume}
\clearpage
\figone{fig9}{Fig.~9 --- fibre orientation}

% =========================================================================
\clearpage
\section*{Formulation audit: equations that changed}
A line-by-line audit against genet23 (Eqs.~1--70, Table~1, App.~A) found the kinematics
(Eqs.~1--16, 21), the fibre field (Eq.~22), the passive, viscous and active laws
(Eqs.~25--33), the pressure and lid loads (Eqs.~34--43, integrated in Eq.~63), the valve law
(Eqs.~35, 36a) and every Table~1 value to be implemented term by term. The items below are
every place where the implementation differed from the paper's \emph{continuous} model
(\S A, now removed), followed by the places where the paper contradicts itself (\S B, kept as
switchable diagnostics whose defaults reproduce the continuous model). Notation:
$\zeta=[\rho,\beta,\varphi,\varepsilon,\eta]$, $r=R+\rho$, $\omega=\dot\beta Z+\dot\varphi$,
$\mathrm d\Omega=R\,\mathrm dR\,\mathrm d\Theta\,\mathrm dZ$.

\subsection*{A. Differences in our model (removed)}

\paragraph{A1. Two-stage Windkessel (Eq.~36b--c).} The paper's circulation is the
C--R--C--R ladder
\begin{equation*}
C_\text{ar}\dot P_\text{ar} = Q_\text{ao} - \frac{P_\text{ar}-P_d}{R_p},\qquad
C_d\dot P_d = \frac{P_\text{ar}-P_d}{R_p} - \frac{P_d-P_\text{vs}}{R_d},\qquad
Q_\text{ao}=K_\text{ar}\langle P_v-P_\text{ar}\rangle_+ .
\end{equation*}
\emph{Old:} two \texttt{BloodVesselRC} blocks. That element is flow-through
($Q_\text{in}=Q_\text{out}$; its capacitor ``stores pressure, not volume''), so the chain solved
\begin{equation*}
C_\text{ar}\dot P_\text{ar} = Q_\text{ao} - \frac{P_\text{ar}-P_d}{R_p},\qquad
C_d\dot P_d = \boxed{Q_\text{ao}} - \frac{P_d-P_\text{vs}}{R_d},
\end{equation*}
i.e.\ the distal node received the full valve flow instead of
$(P_\text{ar}-P_d)/R_p=Q_\text{ao}-C_\text{ar}\dot P_\text{ar}$; $C_\text{ar}$ shaped the
pressure but stored no volume. \emph{New:} svZeroD \texttt{BloodVessel} blocks, which with
$L=0$ solve $P_\text{in}-P_\text{out}=RQ_\text{in}$, $Q_\text{in}-Q_\text{out}=C\dot P_\text{out}$
(capacitor at the outlet): \texttt{aortic\_root}$(R{=}0,C{=}C_\text{ar})\to$
\texttt{arterial}$(R_p,C_d)\to$ \texttt{distal}$(R_d,0)\to P_\text{vs}$, which is Eq.~36b--c
exactly. Effect: ESV $-0.08$\,mL, peak $P$ $-0.01$\,kPa, torsion $+0.01^\circ$.

\paragraph{A2. Upstream element.} \emph{Old:} a \texttt{BloodVessel} between the atrial pressure
and the mitral valve, $P_\text{at}-P_m=R_uQ_\text{mi}$, $Q_\text{mi}-Q_m=C_u\dot P_m$ with
$R_u=\SI{1000}{Pa\,s\,m^{-3}}$ ($0.9\%$ of $K_\text{at}^{-1}$) and $C_u=10^{-12}$\,m$^3$/Pa.
\emph{New:} $R_u=C_u=0$, so the valve sees $P_\text{at}$ as in Eqs.~35/36a. Effect $<0.01^\circ$.

\paragraph{A3. Rigid-body constraints (\S2.4).} The paper imposes $\varphi(R_i)=\eta(R_i)=0$
``during the static simulations'' and states that ``these constraints are not necessary in the
dynamic setting''. \emph{Old:} both pins always on; they replaced the $\hat\varphi(R_i)$ and
$\hat\eta(R_i)$ equations of the virtual-work statement Eq.~45. \emph{New:} pinned only when
$\varrho_0=0$; with inertia all $\hat\varphi,\hat\eta$ equations of Eq.~45 are kept and the rigid
modes follow from momentum conservation (the inner-surface load of Eq.~34 is radial in the
deformed basis, $J\boldsymbol F^{-\top}\boldsymbol e_R=(1{+}\tfrac\rho R)(1{+}\varepsilon)\boldsymbol e_r$,
and the two lid forces of Eq.~43 cancel):
\begin{equation*}
\frac{\mathrm d}{\mathrm dt}\int_\Omega\varrho_0\,r^2\,(\dot\beta Z+\dot\varphi)\,\mathrm d\Omega=0,
\qquad
\frac{\mathrm d}{\mathrm dt}\int_\Omega\varrho_0\,(Z\dot\varepsilon+\dot\eta)\,\mathrm d\Omega=0 .
\end{equation*}
The body now counter-rotates so that the momentum-weighted mean rotation vanishes
($\varphi\approx-\beta L/2$, the mid-height slice stays still), and the twist inertia is taken
about the centroid. For a rigidly moving $\varphi$ field this is the Schur complement of the
mass blocks of Eq.~A8,
\begin{equation*}
M_{\beta\beta}^\text{eff}=M_{\beta\beta}-\frac{M_{\beta\varphi}^2}{M_{\varphi\varphi}}
=2\pi\varrho_0\!\int r^2\Big(\frac{L^3}{3}-\frac{L^3}{4}\Big)R\,\mathrm dR
=2\pi\varrho_0\,\frac{L^3}{12}\!\int r^2R\,\mathrm dR ,
\end{equation*}
one quarter of the pinned value. Effect: tw@550 $7.81\to5.83^\circ$ (the pins had been adding
twist inertia the paper does not have). Kernel switch \texttt{rigid\_pins} ($-1$ auto, 0 free,
1 pinned).

\paragraph{A4. Convective inertia (Eqs.~8, 10, 18).} With Eq.~10,
$\boldsymbol a_c=\mathrm D^2\boldsymbol u(\dot\zeta,\dot\zeta)=2\dot\rho\,\omega\,\boldsymbol e_\theta-r\,\omega^2\boldsymbol e_r$,
and contracting with Eq.~9,
$\hat{\boldsymbol v}=\hat\rho\,\boldsymbol e_r+(\hat\beta Z+\hat\varphi)\,r\,\boldsymbol e_\theta+(Z\hat\varepsilon+\hat\eta)\,\boldsymbol e_z$,
\begin{equation*}
\mathcal P_c(\hat\zeta)=\int_\Omega\varrho_0\,\boldsymbol a_c\cdot\hat{\boldsymbol v}\,\mathrm d\Omega
=2\pi\varrho_0\int_{R_i}^{R_e}\boxed{r}\,\Big[2\dot\rho\,\big(\hat\beta\,I_1+\hat\varphi\,I_0\big)-\hat\rho\,I_2\Big]R\,\mathrm dR,
\end{equation*}
with $I_0=\int_0^L\omega\,\mathrm dZ=\dot\beta\frac{L^2}{2}+\dot\varphi L$,
$I_1=\int_0^L\omega Z\,\mathrm dZ=\dot\beta\frac{L^3}{3}+\dot\varphi\frac{L^2}{2}$,
$I_2=\int_0^L\omega^2\,\mathrm dZ=\dot\beta^2\frac{L^3}{3}+\dot\beta\dot\varphi L^2+\dot\varphi^2L$.
\emph{Old:} the boxed current radius $r$ was the reference radius $R$. \emph{New:} $r=R+\rho$,
with its $\partial/\partial\rho$ in the tangent.

\paragraph{A5. Consistent mass in the current configuration (Eq.~18, A5, A8; commit 027f2ac5).}
$\mathrm D\boldsymbol u^\top\mathrm D\boldsymbol u$ is rotation invariant, so after the $\Theta,Z$
integration only the twist/azimuthal moment arms depend on $\zeta$:
\begin{equation*}
\mathcal M=2\pi\varrho_0\!\int\!\begin{pmatrix}
L&0&0&0&0\\ 0&\tfrac{L^3}{3}\boxed{r^2}&\tfrac{L^2}{2}\boxed{r^2}&0&0\\
0&\tfrac{L^2}{2}\boxed{r^2}&L\,\boxed{r^2}&0&0\\ 0&0&0&\tfrac{L^3}{3}&\tfrac{L^2}{2}\\
0&0&0&\tfrac{L^2}{2}&L\end{pmatrix}R\,\mathrm dR .
\end{equation*}
\emph{Old:} $R^2$ in the boxed entries. \emph{New:} $r^2$; the reference part stays in the
constant mass matrix and $(r^2-R^2)$ is added as a state-dependent force. The Jacobian omits
$\partial\mathcal M/\partial\rho$ (a non-symmetric $\rho$--momentum block that stalls the stiff
test cases); the residual, hence the solution, is exact. Effect $+0.12^\circ$.

\paragraph{A6. Geometry (Table~1).} Table~1 defines the ventricle by $V=\SI{65}{mL}$,
$V_w=\SI{130}{mL}$ and $L/2R_i=1.5$; its radii and length are rounded. \emph{Old:} the rounded
$R_i,R_e,L=19.0,33.0,\SI{57.1}{mm}$ (giving 64.76\,mL, 130.6\,mL, thickness 14.0\,mm).
\emph{New:}
\begin{equation*}
R_i=\Big(\frac{V}{2\pi\cdot1.5}\Big)^{1/3}=\SI{19.035}{mm},\qquad L=3R_i=\SI{57.104}{mm},\qquad
R_e=\sqrt{R_i^2+\frac{V_w}{\pi L}}=\SI{32.969}{mm}.
\end{equation*}
The Fig.~6--8 sweeps use the same construction from 65/130\,mL. Effect: EDV $+0.5$, ESV
$+0.7$\,mL, $-0.1^\circ$.

\paragraph{A7--A9. Initial state, time step, force--velocity smoothing.} The initial cavity
volume is now $V(0)=\pi R_i^2L$ (Eq.~4 at $\rho=\varepsilon=0$) instead of \SI{65}{mL};
$\Delta t=T/800=\SI{1}{ms}$ instead of $T/799$; and $|\dot e_c|$ in Eq.~33 is evaluated as
$\sqrt{\dot e_c^2+\epsilon}$ with $\epsilon=10^{-10}$ instead of $10^{-4}\,$s$^{-2}$. All three
change the results by $<0.02^\circ$ and $<0.1$\,mL.

\paragraph{Net effect (Fig.~5).}
\begin{center}\begin{tabular}{lcccc}
\toprule & old model & new, $n_e{=}3$ & new, $n_e{=}12$ & genet23 \\ \midrule
EDV [mL] & 135.96 & 136.46 & 135.73 & 135.5 \\
ESV [mL] & 74.06 & 74.50 & 73.99 & 74.0 \\
peak $P$ [kPa] & 12.853 & 12.831 & 12.795 & 12.80 \\
torsion at 550\,ms [$^\circ$] & 7.88 & 5.83 & 5.80 & 13.3 \\ \bottomrule
\end{tabular}\end{center}
Pressure and volume now match the paper to $\sim$0.2\,mL; the torsion gap widened because of
A3.

\subsection*{B. Where the paper contradicts itself (diagnostic switches)}

\paragraph{B1. Z-independent velocity memory (App.~A6, A7, A9; kernel \texttt{vmem\_dt}).}
The discrete inertia of the midpoint scheme (Eqs.~67--68) is
\begin{equation*}
\bar{\mathcal P}_a(\hat\zeta)=\int_\Omega2\varrho_0\Big(\mathrm D\boldsymbol u[\zeta^{n+\frac12}]\frac{\zeta^{n+1}-\zeta^n}{\Delta t^2}-\frac{\boldsymbol v^n}{\Delta t}\Big)\cdot\mathrm D\boldsymbol u[\zeta^{n+\frac12}](\hat\zeta)\,\mathrm d\Omega,
\qquad
\boldsymbol v^{n+1}=2\,\mathrm D\boldsymbol u[\zeta^{n+\frac12}]\frac{\zeta^{n+1}-\zeta^n}{\Delta t}-\boldsymbol v^n .
\end{equation*}
Eq.~A7 integrates the first term to $\mathcal M$ (Eq.~A8) but writes the second as
$-\mathcal N^\top\boldsymbol v^n/\Delta t$ with $\mathcal N=\int_0^L\mathrm D\boldsymbol u\,\mathrm dZ$
(Eq.~A9), i.e.\ $\boldsymbol v^n=\boldsymbol v^n(R)$ is taken out of the $Z$-integral, and \S2.7.2
stores $\boldsymbol v$ ``at the quadrature point level, similarly to internal variables''. The
exact velocity $\mathrm D\boldsymbol u\,\dot\zeta$ depends on $Z$
($v_\theta=(\dot\beta Z+\dot\varphi)r$, $v_z=Z\dot\varepsilon+\dot\eta$). Consistent storage would
give $\mathcal N^\top\boldsymbol v^n=\mathcal M\dot\zeta^n$; storing the $Z$-mean
$\bar{\boldsymbol v}=\mathcal N\dot\zeta/L$ keeps only $(\mathcal N^\top\mathcal N/L)\,\dot\zeta$
and drops the momentum $\Delta\mathcal M\,\dot\zeta$ every step, with
\begin{equation*}
\Delta\mathcal M=\int_0^L\mathrm D\boldsymbol u^\top\mathrm D\boldsymbol u\,\mathrm dZ
-\frac1L\Big(\int_0^L\mathrm D\boldsymbol u\,\mathrm dZ\Big)^{\!\top}\Big(\int_0^L\mathrm D\boldsymbol u\,\mathrm dZ\Big)\;\succeq0,
\end{equation*}
with, at small twist, $\Delta\mathcal M_{\beta\beta}=r^2L^3/12$,
$\Delta\mathcal M_{\varepsilon\varepsilon}=L^3/12$ and all other entries $O((\beta L)^2)$.
For a mode of mass $m$ whose memory is discarded,
$2m(q^{n+1}-q^n)/\Delta t^2=F^{n+\frac12}$, i.e.\ $(q^{n+1}-q^n)/\Delta t=F\,\Delta t/(2m)$: a
dashpot $2m/\Delta t$. The emulation therefore adds the generalized force
\begin{equation*}
f_\text{vm}(\hat\zeta)=\frac{2}{\Delta t_p}\,2\pi\!\int\varrho_0\,\hat\zeta^\top\Delta\mathcal M\,\dot\zeta\;R\,\mathrm dR ,
\end{equation*}
with $\Delta t_p$ the paper's step (\texttt{vmem\_dt}; $\Delta\mathcal M$ from an 8-point Gauss rule
in $Z$ with the exact $\mathrm D\boldsymbol u$ of Eq.~A5). It damps only the $Z$-deviation of the
twist and axial-stretch velocities, grows as $1/\Delta t_p$, and scales with $L^3$ (physical
twist viscosity scales with $L$). At $\Delta t_p=\SI{1}{ms}$ it is $\approx15\times$ the
tissue-viscosity twist damping $\gamma\int\tfrac{R^2}{2}\mathrm d\Omega$.

\paragraph{B2. Active stress, Eq.~59 vs.\ Eq.~30 (kernel \texttt{active\_stretch\_pow}).}
\begin{equation*}
\sigma_{1D}^{n+\frac12\#}=\frac{T_\text{fib}^{n+\frac12\#}}{1+2\boldsymbol e_F\!\cdot\!\boldsymbol E^n\!\cdot\!\boldsymbol e_F}\,
\frac{\boldsymbol e_F\cdot(\boldsymbol E^{n+1}-\boldsymbol E^n)\cdot\boldsymbol e_F}{\boldsymbol e_F\cdot\Delta t\,\dot{\boldsymbol E}^{n+\frac12\#}\cdot\boldsymbol e_F}
\;\xrightarrow{\;\Delta t\to0\;}\;\frac{T_\text{fib}}{(1+e_\text{fib})^2},
\end{equation*}
whereas Eq.~30 gives $\sigma_{1D}=T_\text{fib}/(1+e_\text{fib})$,
since $1+2\boldsymbol e_F\cdot\boldsymbol E\cdot\boldsymbol e_F=I_4=(1+e_\text{fib})^2$. The energy
identity of Eq.~46,
$\Sigma^a:\dot{\boldsymbol E}=\sigma_{1D}(1+e_\text{fib})\dot e_\text{fib}=T_\text{fib}\dot e_\text{fib}$,
requires Eq.~30. Switch: $\sigma_{1D}=T_\text{fib}/(1+e_\text{fib})^p$, $p=1$ (default) or 2.

\paragraph{B3. Missing cylindrical metric (Eqs.~A7, A10, A11; kernel \texttt{inertia\_metric}).}
All three are printed as $\int\ldots\mathrm dR$ instead of $\int\ldots R\,\mathrm dR$. For the
stiffness and incompressibility terms this cannot be the code (the pressure--volume balance
would change by $\sim1/\bar R$, and ours matches with the Table~1 parameters); the switch drops
$R$ from the inertia terms only.

\paragraph{B4. Internal-variable scheme of Ref.~6 (Chapelle 2012, Eq.~26; kernel \texttt{tauc\_alpha}).}
genet23 integrates $(k_c,\,\tau_c/\sqrt{k_c})$ with the scheme of Ref.~6. There, the $k_c$ update
keeps $\alpha$, but the printed $X=\tau_c/\sqrt{k_c}$ update has the decay
$-\tfrac12\big(|u^{n+\frac12}|+|\tfrac{e_c^{n+1}-e_c^n}{\Delta t}|\big)X^{n+1}$, without $\alpha$.
With $\tau_c=\sqrt{k_c}\,X$ the source terms cancel and the continuous limit is
\begin{equation*}
\dot\tau_c=-\Big(|\nu|+\tfrac{\alpha+1}{2}|\dot e_c|\Big)\tau_c+k_c\dot e_c+n_0\sigma_0|\nu|_+ ,
\qquad
\dot k_c=-\big(|\nu|+\alpha|\dot e_c|\big)k_c+n_0k_0|\nu|_+ ,
\end{equation*}
i.e.\ Eq.~33 with $\alpha_\tau=(\alpha+1)/2=6.5$ in the $\tau_c$ equation only (consistency with
Eq.~33 requires $-\tfrac12(|u|+\alpha|\dot e_c|)X$). Switch: \texttt{tauc\_alpha} ($-1$: same as
$\alpha$).

\paragraph{Results (Fig.~5).}
\begin{center}\small\begin{tabular}{lccccc}
\toprule & torsion RMS [$^\circ$] & volume RMS [mL] & ESV [mL] & peak $P$ [kPa] & tw@550 [$^\circ$]\\ \midrule
faithful model (\S A) & 3.58 & 2.72 & 74.50 & 12.83 & 5.83 \\
B1, $\Delta t_p=\SI{1}{ms}$ (paper) & 4.16 & 1.28 & 73.16 & 12.89 & 15.47 \\
B1, $\Delta t_p=\SI{2}{ms}$ (half strength) & 1.82 & 1.28 & 74.15 & 12.89 & 14.41 \\
B2, $p=2$ & 5.12 & 5.07 & 79.59 & 11.68 & 5.40 \\
B3, SI units & 3.88 & 6.07 & 76.70 & 12.85 & 16.14 \\
B4, $\alpha_\tau=6.5$ & 3.39 & 2.75 & 69.41 & 13.85 & 7.36 \\ \midrule
genet23 & --- & --- & 74.0 & 12.80 & 13.3 \\ \bottomrule
\end{tabular}\end{center}
B2, B3 and B4 break pressure and volume and are ruled out; B4 also leaves the untwist tail
and the aspect-ratio dependence unchanged (Fig.~6 slope $-0.3^\circ$), since the $\alpha|\dot e_c|$
term acts during shortening while the diastolic decay is dominated by $|\nu|$. B1 changes only the twist and
reproduces the aspect-ratio dependence of the residual twist at 600\,ms (Fig.~6,
aspect $1.0\to2.0$: $+10.4$ to $+10.9^\circ$ vs.\ the paper's $+10.6^\circ$; faithful model
$-0.9^\circ$), but at the paper's $\Delta t=\SI{1}{ms}$ it is about twice too strong, and the
convergence study below shows it is incompatible with the paper's own Fig.~4.

\begin{center}\includegraphics[width=0.92\linewidth]{audit_fig5.png}\end{center}

% =========================================================================
\clearpage
\section*{Convergence of the faithful model and of B1}
Fig.~5 baseline, last of 12 cycles, compared with the finest run of each study
($\Delta t=\SI{0.125}{ms}$, resp.\ $n_e=12$) using the paper's metrics, Eqs.~69--70 (normalized
RMS errors of $P$ and $V$, relative error of the PV work $W=\sum P^{n+\frac12}(V^{n+1}-V^n)$), and
the same RMS metric for the torsion. Three temporal studies: (a) the faithful model; (b) B1 with
the paper's step held at $\Delta t_p=\SI{1}{ms}$ while our step is refined (is the emulated
effect resolved by our solver?); (c) B1 with $\Delta t_p$ equal to the step --- what
CardiacLab's own time-step study would give if its scheme dropped the $Z$-deviation momentum.
\begin{center}\small\setlength{\tabcolsep}{4pt}\begin{tabular}{r|ccc|ccc|cccc|c}
\toprule
 & \multicolumn{3}{c|}{(a) faithful} & \multicolumn{3}{c|}{(b) B1, $\Delta t_p{=}1$\,ms} &
 \multicolumn{4}{c|}{(c) B1, $\Delta t_p{=}\Delta t$} & genet23 \\
$\Delta t$ [ms] & $e_{PV}$ & $e_\text{tors}$ & tw@550 & $e_{PV}$ & $e_\text{tors}$ & tw@550 &
 $e_{PV}$ & $e_\text{tors}$ & tw@550 & EDV & Fig.~4 $e_{PV}$ \\ \midrule
2    & 3.0e-4 & 8.6e-4 & 5.82 & 2.8e-4 & 3.1e-4 & 15.47 & 3.9e-2 & 5.1e-1 & 14.40 & 136.1 & $\approx$2.4e-3 \\
1    & 6.6e-5 & 2.3e-4 & 5.83 & 6.2e-5 & 6.7e-5 & 15.47 & 3.7e-2 & 3.7e-1 & 15.47 & 135.2 & $\approx$6e-4 \\
0.5  & 1.6e-5 & 5.9e-5 & 5.83 & 1.8e-5 & 2.1e-5 & 15.47 & 2.7e-2 & 2.1e-1 & 15.97 & 133.3 & $\approx$1.8e-4 \\
0.25 & 4.2e-6 & 1.2e-5 & 5.83 & 2.8e-6 & 3.8e-6 & 15.47 & 1.1e-2 & 8.4e-2 & 16.16 & 131.4 & $\approx$4.5e-5 \\
0.125 & \multicolumn{2}{c}{ref.} & 5.83 & \multicolumn{2}{c}{ref.} & 15.47 & \multicolumn{2}{c}{ref.} & 16.18 & 130.2 & ref. \\
\bottomrule
\end{tabular}\end{center}
{\footnotesize genet23 Fig.~4 values are read off the published plot (approximate).}
\begin{itemize}
\item \textbf{(a)} The faithful model converges at second order in time; torsion is
  $\Delta t$-independent to $0.01^\circ$. Its PV-work errors are about $10\times$ below the
  paper's with the same slope (different scheme, same behaviour).
\item \textbf{(b)} B1 at fixed $\Delta t_p$ is resolved by our solver just as well, so the
  $+9.6^\circ$ it adds at 550\,ms is not an artifact of our discretization.
\item \textbf{(c)} If CardiacLab dropped the $Z$-deviation momentum, its time-step study would
  show PV-work errors of 1--4\% that barely decrease, an end-diastolic volume drifting by
  \SI{5.9}{mL} between 2 and \SI{0.125}{ms}, and a torsion that never converges. The
  published Fig.~4 shows $\sim$0.06\% at \SI{1}{ms} and clean second-order convergence, as in (a).
  \textbf{The paper's own convergence study therefore rules out B1 as written.} The PV-work
  error in (c) comes mainly from the \emph{twist} part of $\Delta\mathcal M$, so a twist-only
  loss does not escape Fig.~4 either (see ``Scaling fingerprint'' below).
\end{itemize}
\begin{center}\begin{tabular}{r|cccc|cccc}
\toprule
 & \multicolumn{4}{c|}{faithful} & \multicolumn{4}{c}{B1, $\Delta t_p{=}1$\,ms} \\
$n_e$ & $e_{PV}$ & $e_\text{tors}$ & tw@550 & EDV & $e_{PV}$ & $e_\text{tors}$ & tw@550 & EDV \\ \midrule
2  & 1.1e-2 & 1.6e-2 & 5.87 & 137.13 & 1.5e-2 & 8.5e-3 & 15.53 & 135.78 \\
3  & 6.9e-3 & 6.3e-3 & 5.83 & 136.46 & 6.8e-3 & 8.5e-3 & 15.47 & 135.16 \\
6  & 1.5e-3 & 1.3e-3 & 5.80 & 135.90 & 1.5e-3 & 1.7e-3 & 15.44 & 134.67 \\
12 & \multicolumn{2}{c}{ref.} & 5.80 & 135.73 & \multicolumn{2}{c}{ref.} & 15.42 & 134.52 \\
\bottomrule
\end{tabular}\end{center}
Spatially both converge at second order (P1 elements); torsion is mesh-independent to $0.05^\circ$
from $n_e=3$, and $n_e=6$ is within \SI{0.2}{mL} of $n_e=12$ (the figures above use $n_e=6$).
The paper's single P3 element reaches $\sim$3e-5 (Fig.~4, upper left), which our P1 mesh would
need $n_e\gtrsim24$ for; $n_e=24$ does not converge at present (open issue).

\begin{center}\includegraphics[width=\linewidth]{convergence.png}\end{center}

% =========================================================================
\clearpage
\section*{Literature formulation screen (INRIA/MEDISIM and colleagues)}
Alternative formulations of each model part, taken from the authors' group and colleagues
(papers in the Zotero library), were swapped one at a time into the faithful model.
Each one is a kernel switch whose default reproduces genet23 bit-for-bit.
All runs use the Fig.~5 set-up ($n_e=3$, \SI{1}{ms}, 12 cycles, last cycle) and are scored
against the digitised Figs.~5, 6 and 9. The best rows were re-run at $n_e=6$, with changes below
$0.05^\circ$.

\paragraph{Formulations screened.}
\begin{center}\footnotesize\begin{tabular}{P{2.1cm}P{8.9cm}P{4.6cm}}
\toprule
part & alternatives (source) & kernel switch \\ \midrule
series law & additive in engineering strain, $e_\text{fib}=e_s+e_c$ (genet23); multiplicative in
Green--Lagrange strain, $(1{+}2e_\text{1D})=(1{+}2e_s)(1{+}2e_c)$,
$\sigma_\text{1D}=\frac{1+2e_c}{1+2e_\text{1D}}T$ (Sainte-Marie 2006, Chapelle 2012, Caruel 2013)
& \texttt{series\_law} \\
Frank--Starling & $n_0(e_c)$ on the BCS sources (genet23); $d(e_c)$ on the output tension
(Sainte-Marie 2006) & \texttt{starling\_out} \\
BCS kinetics & $\alpha=1$ (Bestel 2001); $\alpha=1.5$, $k_0/\sigma_0=0.8$ (Caruel 2013);
$\sigma_0=(\frac12+s_0)k_0$ at $s_0=0$ (Chapelle 2012); $\sigma_0/k_0=2$ (Sainte-Marie 2006)
& \texttt{bcs\_alpha}, \texttt{k\_0} \\
activation $\nu(t)$ & diastolic value $-40$ (Kimmig), $-100$ (Ruijsink 2020); $-13$ with a
\SI{120}{ms} upstroke (Caruel 2013) & activation function \\
activation timing & transmural delay endo$\to$epi or epi$\to$endo (Chabiniok 2012) &
\texttt{act\_delay} \\
passive & Eq.~25 with the pure-exponential sets of Chapelle/Le~Gall 2023, Le~Gall 2020 and
Caruel 2013; isotropic Mooney--Rivlin (Chabiniok 2012); Fung
$W=\frac C2(e^Q-1)$ in the fibre basis $(\mathbf e_F,\mathbf e_R,\mathbf e_F{\times}\mathbf e_R)$
(Guccione 1991, Genet 2014, Lee 2014, Xi 2011); Holzapfel--Ogden (Pfaller 2019 with sheet in-plane or
transmural, Asner 2015, Hadjicharalambous 2017, Nordsletten 2021) &
\texttt{passive\_law}, \texttt{fung\_*}, \texttt{ho\_*} \\
viscous & $\gamma\dot{\mathbf E}$ (genet23; Pfaller 2019 with $\eta=\SI{100}{Pa\,s}$);
$\gamma\,\mathrm{tr}(\dot{\mathbf E})\mathbf I$ (Chapelle 2012 as printed); anisotropic
Kelvin--Voigt $\Sigma^v_{ab}=\gamma_{ab}\dot E_{ab}$ in the fibre basis, the \SI{10}{rad/s}
surrogate of the fractional law of Nordsletten 2021 ($\gamma_{FF,SS,NN,FS,FN,SN}=143, 78, 36, 269,
153, \SI{126}{Pa\,s}$) & \texttt{visc\_mode}, \texttt{gam\_*} \\
boundary & epicardial normal spring--dashpot (Pfaller 2019 Eq.~5); omnidirectional Robin
(Moireau 2011 Eq.~9); epicardial tangential dashpot; tangential visco-elastic support of the
end annuli (Chabiniok 2012 Eq.~13, Sainte-Marie 2006 base) & \texttt{epi\_*},
\texttt{support\_*} \\ \bottomrule
\end{tabular}\end{center}
The passive laws are fitted to other species and data sets, and taken literally they move the EDV
by $-60$ to $+160$\,mL. So each one is shown after scaling its moduli (not its exponents) by
one factor that restores the faithful EDV of \SI{136.5}{mL}. Taken literally, all of them are
worse still.
Caruel's length-dependent relaxation was not implemented, since the kinetics class as a whole
already fails.

\paragraph{Scorecard.} RMS errors against the digitised curves; tw@550 and tw@650 in degrees;
``slope'' is the Fig.~6 torsion at \SI{600}{ms} for aspect ratio 2.0 minus 1.0.
\begin{center}\footnotesize\setlength{\tabcolsep}{3.5pt}\begin{longtable}{lcccccccc}
\toprule
 & \multicolumn{3}{c}{Fig.~5 RMS} & & & Fig.~6 & Fig.~6 & Fig.~9 \\
variant & tors.\ [$^\circ$] & vol.\ [mL] & $P$ [kPa] & tw@550 & tw@650 & slope & RMS & RMS \\ \midrule
\textbf{faithful model} & 3.58 & 2.72 & 0.097 & 5.83 & $-2.23$ & $-0.8$ & 3.42 & 3.23 \\
\textbf{genet23} & --- & --- & --- & 13.3 & 3.8 & $+10.5$ & --- & --- \\ \midrule
\multicolumn{9}{l}{\emph{active law and activation}} \\
GL series law (Chapelle 2012) & 3.82 & 2.82 & 0.500 & 5.98 & $-2.14$ & $-0.6$ & 3.62 & 3.19 \\
Starling on output (Sainte-Marie 2006) & 3.42 & 2.63 & 0.156 & 6.88 & $-2.13$ & $-0.5$ & 3.25 & 2.91 \\
$\alpha=1$ (Bestel 2001) & 3.09 & 5.87 & 1.250 & 9.24 & $-1.86$ & $+0.6$ & 3.09 & 2.88 \\
$\alpha=1.5$ (Caruel 2013) & 3.10 & 5.46 & 1.175 & 9.08 & $-1.89$ & $+0.5$ & 3.09 & 2.89 \\
$k_0=\SI{130}{kPa}$ (Chapelle 2012) & 3.92 & 2.80 & 0.221 & 5.03 & $-2.25$ & $-0.9$ & 3.78 & 3.49 \\
$k_0=\SI{52}{kPa}$ (Caruel 2013) & 4.32 & 3.03 & 0.318 & 4.06 & $-2.27$ & $-1.1$ & 4.18 & 3.73 \\
$k_0=\SI{32.5}{kPa}$ (Sainte-Marie 2006) & 4.47 & 3.12 & 0.344 & 3.71 & $-2.28$ & $-1.1$ & 4.32 & 3.81 \\
$\nu$ diastolic $-40$ (Kimmig) & 4.07 & 4.57 & 0.180 & 3.50 & $-2.31$ & $-1.2$ & 3.85 & 3.54 \\
$\nu$ $-13$, \SI{120}{ms} ramp (Caruel 2013) & 5.85 & 21.02 & 2.248 & 20.11 & 17.35 & $+12.7$ & 6.32 & 4.33 \\
$\nu$ diastolic $-100$ (Ruijsink 2020) & 6.15 & 12.32 & 0.775 & $-1.54$ & $-2.42$ & $-1.6$ & 5.88 & 4.97 \\
delay endo$\to$epi \SI{20}{ms} (Chabiniok 2012) & 6.21 & 3.24 & 0.538 & 7.42 & $-2.18$ & $-1.0$ & 5.98 & 5.26 \\
delay epi$\to$endo \SI{20}{ms} & 8.29 & 4.47 & 0.484 & $-0.26$ & $-2.35$ & $-1.4$ & 8.22 & 7.23 \\ \midrule
\multicolumn{9}{l}{\emph{passive law (EDV-matched modulus scale)}} \\
Chapelle/Le~Gall 2023 ($\times0.42$) & 4.94 & 3.13 & 0.190 & 4.21 & $-4.38$ & $-2.8$ & 4.68 & 4.71 \\
Le~Gall 2020 ($\times0.57$) & 4.67 & 2.81 & 0.165 & 4.03 & $-3.87$ & $-2.4$ & 4.41 & 4.40 \\
Caruel 2013 ($\times4.8$) & 4.53 & 2.66 & 0.084 & 1.74 & $-0.51$ & $-0.3$ & 4.37 & 3.80 \\
isotropic MR, Chabiniok 2012 ($\times0.28$) & 5.52 & 12.08 & 0.154 & 0.67 & 0.02 & $0.0$ & 5.35 & 4.43 \\
Fung, Genet 2014 ($\times1.43$) & 4.67 & 2.60 & 0.150 & 2.39 & $-2.86$ & $-1.8$ & 4.38 & 3.93 \\
Fung, Guccione 1991, $\pm60^\circ$ ($\times0.17$) & 5.54 & 3.26 & 0.189 & 2.60 & $-4.68$ & $-3.0$ & 5.25 & 4.76 \\
Fung, Guccione 1991, $75/{-45}^\circ$ ($\times0.23$) & 18.80 & 2.68 & 0.124 & $-2.64$ & $-6.09$ & $-3.9$ & 18.50 & n/a \\
Fung, Lee 2014 ($\times0.22$) & 5.49 & 4.50 & 0.197 & 6.09 & $-5.88$ & $-3.4$ & 5.28 & 5.14 \\
Fung, Xi 2011 ($\times0.082$) & 5.67 & 5.46 & 0.210 & 8.27 & $-6.65$ & $-3.5$ & 5.52 & 5.28 \\
HO, Pfaller 2019, sheet in-plane ($\times0.011$) & 6.69 & 5.73 & 0.363 & 4.65 & $-5.76$ & $-3.7$ & 6.38 & 6.71 \\
HO, Pfaller 2019, sheet transmural ($\times0.018$) & 8.18 & 4.92 & 0.414 & 3.47 & $-7.53$ & $-5.0$ & 7.89 & 9.29 \\
reduced HO, Hadjicharalambous 2017 ($\times0.69$) & 4.98 & 1.87 & 0.196 & 3.13 & $-4.03$ & $-2.6$ & 4.70 & 4.76 \\
reduced HO, Asner 2015 ($\times0.35$) & 5.44 & 1.49 & 0.235 & 2.03 & $-4.19$ & $-2.8$ & 5.14 & 5.43 \\
HO, Nordsletten 2021 ($\times0.010$) & 9.46 & 6.31 & 0.475 & 4.16 & $-8.87$ & $-5.8$ & 9.17 & 10.20 \\ \midrule
\multicolumn{9}{l}{\emph{viscous law}} \\
$\gamma\,\mathrm{tr}(\dot{\mathbf E})\mathbf I$ (Chapelle 2012 as printed) & 5.55 & 5.50 & 0.120 & $-5.45$ & $-2.32$ & $+1.8$ & 5.36 & 4.83 \\
isotropic $\eta=\SI{100}{Pa\,s}$ (Pfaller 2019) & 2.85 & 1.69 & 0.082 & 7.99 & $-1.50$ & $+0.6$ & 2.74 & 2.69 \\
anisotropic KV (Nordsletten 2021), $S$ transmural & 2.67 & 1.43 & 0.078 & 8.34 & $-1.17$ & $+1.0$ & 2.57 & 2.56 \\
anisotropic KV (Nordsletten 2021), $S$ in-plane & 1.74 & 1.11 & 0.065 & 9.97 & 0.54 & $+2.5$ & 1.82 & 1.83 \\ \midrule
\multicolumn{9}{l}{\emph{boundary support, excluded by genet23 (Sec.~2.4, Eq.~44), diagnostic only; $k$ in Pa/m, $c$ in Pa\,s/m}} \\
epicardium normal, $k=2\text{e}5$, $c=5\text{e}3$ (Pfaller 2019) & 2.98 & 9.64 & 0.238 & 7.21 & $-0.60$ & $+0.2$ & 2.97 & 2.54 \\
epicardium omni, $k=1\text{e}4$, $c=1\text{e}5$ (Moireau 2011) & 11.02 & 35.10 & 0.938 & 17.07 & 16.71 & $+10.9$ & 11.00 & 9.17 \\
epicardium circumferential, $c=1\text{e}3$ & 2.42 & 2.32 & 0.089 & 9.55 & $-1.00$ & $+4.4$ & 2.21 & 2.33 \\
epicardium circumferential, $c=5\text{e}3$ & 0.82 & 1.35 & 0.093 & 14.62 & 5.17 & $+11.8$ & 1.41 & 1.18 \\
\quad + axial, $c=5\text{e}3$ & 0.82 & 1.35 & 0.088 & 14.68 & 5.19 & $+11.9$ & 1.41 & 1.15 \\
\quad + normal, $c=5\text{e}3$ & 1.27 & 1.58 & 0.064 & 15.15 & 6.28 & $+12.4$ & 1.85 & 1.35 \\
end annuli, $c=500$ & 3.07 & 2.55 & 0.093 & 7.63 & $-1.81$ & $+0.7$ & 2.92 & 2.81 \\
end annuli, $c=2000$ & 1.76 & 2.11 & 0.087 & 11.00 & 0.17 & $+4.3$ & 1.77 & 1.87 \\
end annuli, $c=4000$ & 0.45 & 1.68 & 0.088 & 13.28 & 2.96 & $+7.0$ & 1.04 & 1.13 \\
end annuli, $c=8000$ & 1.77 & 1.13 & 0.102 & 15.31 & 7.01 & $+9.4$ & 2.06 & 1.68 \\
base annulus only, $c=8000$ & 3.49 & 2.64 & 0.075 & 7.64 & $-3.04$ & $-0.5$ & 3.33 & 3.21 \\
end annuli, $k=2\text{e}5$, $c=5\text{e}3$ & 6.04 & 3.82 & 0.170 & 0.68 & $-0.59$ & $-0.1$ & 5.72 & 4.38 \\
\bottomrule
\end{longtable}\end{center}
{\footnotesize Guccione 1991 fitted its set for $75/{-45}^\circ$ fibres (its $\pm60^\circ$ fit has a
negative shear coefficient); the Fig.~9 fibre sweep overrides those fibres (n/a).}

\paragraph{Findings.}
\begin{itemize}
\item \textbf{No constitutive alternative closes the gap.} Changes to the active law, kinetics or
  activation either leave the untwisting unchanged or buy twist at the expense of pressure and
  volume. Examples: $\alpha=1$ gives $+3.4^\circ$ at 550\,ms for 1.25\,kPa of pressure RMS; the
  Caruel ramp gives a twist that never relaxes and a \SI{21}{mL} volume RMS.
\item \textbf{Every literature passive law makes torsion worse.} All of them carry 20--250$\times$
  the small-strain shear stiffness of Eq.~25 (\SI{14}{Pa}). At equal EDV this extra stiffness
  pulls the twist back sooner and overshoots it negative (tw@650 down to $-9^\circ$).
\item \textbf{Diagnostic: how much twist damping is missing.} The boundary rows are not
  admissible (next item). They serve only as a probe. An added viscous resistance to wall
  rotation that leaves volume alone reproduces Figs.~5, 6, 9 (and 7, 8) together with $P$ and
  $V$. The effect collapses onto one number, the added twist damping $D_\beta$ relative to the
  twist damping of Eq.~28 itself (small strain, $E_{\Theta Z}=R\beta/2$, rigid rotation
  self-centred at mid-length):
  \begin{equation*}
  D_\gamma=\tfrac{\pi}{4}\gamma L(R_e^4-R_i^4),\qquad
  D_\text{ends}=\tfrac{\pi}{4}cL^2(R_e^4-R_i^4)=\tfrac{cL}{\gamma}D_\gamma,\qquad
  D_\text{epi}=\tfrac{\pi}{6}cR_e^3L^3 .
  \end{equation*}
  \begin{center}\small\begin{tabular}{lccccccccc}
  \toprule
  $D_\beta/D_\gamma$ & 0 & 0.41 & 0.43$^*$ & 1.06 & 1.63 & \textbf{3.26} & 5.30 & 6.53 & genet23 \\
  applied as & --- & ends & $\gamma{=}100$ & epi & ends & \textbf{ends} & epi & ends & \\ \midrule
  tw@550 [$^\circ$] & 5.83 & 7.63 & 7.99 & 9.55 & 11.00 & \textbf{13.28} & 14.62 & 15.31 & 13.3 \\
  tw@650 [$^\circ$] & $-2.23$ & $-1.81$ & $-1.50$ & $-1.00$ & 0.17 & \textbf{2.96} & 5.17 & 7.01 & 3.8 \\
  \bottomrule
  \end{tabular}\end{center}
  {\footnotesize $^*$isotropic $\gamma$ also damps the volume modes.}\\
  The paper's curves need $D_\beta\approx3.3\,D_\gamma$, about $1.1\times10^{-5}$\,N\,m$^3$\,s, whether
  it comes from the end annuli ($c\approx\SI{4e3}{Pa\,s/m}$, chosen by interpolating the
  500/2000/8000 sweep to the Fig.~5 tw@550) or the epicardium ($c\approx\SI{3e3}{Pa\,s/m}$).
  The Fig.~5 match is therefore partly by construction; the agreement of Figs.~6--9 and of the
  untwist shape is not tuned. As tissue viscosity it would be an extra
  $\approx\SI{230}{Pa\,s}$ acting on $E_{\Theta Z}$ alone. The anisotropic Nordsletten viscosity
  gets part of the way ($D_\beta/D_\gamma\approx1$) because it weights shear. Placement shows
  only in the aspect-ratio slope: $D_\text{ends}\propto L^2$ gives $+7.0$, $D_\text{epi}\propto L^3$
  gives $+11.8$, against the paper's $+10.5$.
\item \textbf{Boundary support is ruled out by the paper.} Directly after Eq.~44, genet23
  names visco-elastic conditions on the external surface for the pericardial fluid (Refs.~8,
  41) and states that they ``will not be considered in this paper''. Eq.~44 gives the complete
  external virtual power as the endocardial pressure plus the two lid pressure terms. The lids
  couple to the wall only through the axial motion (Eqs.~39--42), so no external force acts on
  the rotation, at the epicardium or at the ends. The magnitude is not a literature value
  either: Pfaller's $c_p=\SI{5e3}{Pa\,s/m}$ is normal-only, and a normal dashpot damps volume,
  not twist.
\item \textbf{Status.} The missing $\approx3.3\,D_\gamma$ of twist-selective damping comes from
  neither the written boundary conditions nor any of the 30 constitutive alternatives. The next
  section narrows down what kind of term it must be.
\end{itemize}
\begin{center}\includegraphics[width=0.88\linewidth]{fig5_ends4e3.png}\\
\includegraphics[width=0.88\linewidth]{fig6_epiT5e3.png}\\
{\footnotesize Diagnostic probes only: both terms are excluded by genet23 (Sec.~2.4, Eq.~44).
The end-annulus $c$ was chosen to match the Fig.~5 tw@550.}\end{center}

% =========================================================================
\section*{Scaling fingerprint of the missing twist damping}
\paragraph{Method.} A pure generalized twist dashpot $D\dot\beta$ on the $\beta$ equation (kernel
\texttt{twist\_c}; power $D\dot\beta^2$, no approximation) is used as a probe. For every case of
Figs.~5--9, $D$ was scanned ($0$--$11\,D_{\gamma,0}$, 13 values) and the value minimising the
torsion RMS against that case's digitised curve was taken. Two further estimators are the RMS
over $t\ge\SI{420}{ms}$ (untwist only) and matching tw@550. The baseline appears in five figures;
its five traces give $D/D_{\gamma,0}=3.6$--$5.0$ (full RMS), so digitisation scatter is about
15\%. Each case is therefore normalised by the baseline trace of its own figure. Each candidate
mechanism has a definite scaling of $D$ with length $L$ and $Q=R_e^4-R_i^4$:
\begin{itemize}
\item any local constitutive term (viscosity, fibre dashpot, stiffness change): $D\propto LQ$, since
  $\beta$ enters every strain as $R\beta$, uniform in $Z$;
\item a boundary term at the ends: $D\propto L^2Q$;
\item $Z$-momentum (inertia, or the B1 momentum loss): $D\propto L^3Q$.
\end{itemize}

\paragraph{Result.}
\begin{center}\small\setlength{\tabcolsep}{4pt}\begin{tabular}{lccccc|ccc}
\toprule
 & & & \multicolumn{3}{c|}{measured $D/D_\text{base}$} & \multicolumn{3}{c}{predicted} \\
case & $L$ [mm] & $Q/Q_0$ & full & $\ge$420 & tw@550 & $LQ$ & $L^2Q$ & $L^3Q$ \\ \midrule
aspect 1.0 & 43.6 & 1.72 & 0.50 & 0.47 & 0.53 & 1.31 & 1.00 & 0.76 \\
aspect 2.0 & 69.2 & 0.68 & 1.56 & 1.55 & 1.49 & 0.83 & 1.00 & 1.21 \\
$V_\text{int}$ 58 & 55.0 & 1.02 & 0.85 & 0.84 & 0.80 & 0.98 & 0.95 & 0.91 \\
$V_\text{int}$ 72 & 59.1 & 0.98 & 1.17 & 1.14 & 1.19 & 1.02 & 1.05 & 1.09 \\
$V_\text{wall}$ 117 & 57.1 & 0.85 & 0.88 & 0.88 & 0.96 & 0.86 & 0.86 & 0.86 \\
$V_\text{wall}$ 143 & 57.1 & 1.15 & 1.11 & 1.11 & 1.10 & 1.15 & 1.15 & 1.16 \\
fibre $30^\circ$ & 57.1 & 1 & 0.78 & 0.73 & 0.75 & 1 & 1 & 1 \\
fibre $45^\circ$ & 57.1 & 1 & 0.95 & 0.92 & 0.87 & 1 & 1 & 1 \\
fibre $75^\circ$ & 57.1 & 1 & 1.06 & 1.04 & 0.95 & 1 & 1 & 1 \\
fibre $90^\circ$ & 57.1 & 1 & 1.03 & 0.94 & 0.60 & 1 & 1 & 1 \\ \bottomrule
\end{tabular}\end{center}
A free fit over Figs.~6--8 gives $D\propto L^{a}Q^{b}$ with $a=4.2$--$4.6$, $b=0.6$--$0.85$
(depending on the estimator; $a=4.3$--$4.6$ with $b=1$ fixed).
\begin{itemize}
\item \textbf{Radial weighting (Fig.~8, $L$ fixed):} $D\propto Q$ fits (0.88/1.11 measured vs.\
  0.86/1.15), so the missing term acts on the circumferential velocity $r\dot\beta$, like every candidate.
\item \textbf{Length (Figs.~6, 7):} the required damping grows steeply with length, from half the
  baseline value at aspect 1.0 to $1.5\times$ at aspect 2.0. A local constitutive term predicts
  the opposite trend ($1.31\to0.83$). \emph{No local material law, written or not, can
  produce the paper's aspect dependence}: it keeps the twist relaxation time independent of
  $L$, whereas the paper's grows with $L$. End-boundary damping gives no trend. Only
  $Z$-momentum terms go the right way, and they reach about half the measured growth
  ($L^3$ vs.\ $L^{4.4}$).
\item \textbf{Fibre angle (Fig.~9, geometry fixed):} $D$ is roughly independent of the angle
  (0.75--1.06; the $30^\circ$ case fits poorly). This is consistent with a non-material term, but
  it is a weak test.
\item \textbf{Start-up and cycle index:} cycle 5 (the paper's) and cycle 12 differ by
  $0.1^\circ$ at 550\,ms (5.93 vs.\ 5.83), because the twist mode forgets its initial state
  within a cycle.
\end{itemize}

\paragraph{Twist-only velocity loss vs.\ Fig.~4 (kernel \texttt{vmem\_mask}).} B1 was restricted
to the in-plane (twist) or the axial velocity, and the paper-style time-step study was repeated
($\Delta t_p=\Delta t$, reference $\Delta t=\SI{0.125}{ms}$):
\begin{center}\small\begin{tabular}{r|ccc|ccc|c}
\toprule
 & \multicolumn{3}{c|}{twist-only loss} & \multicolumn{3}{c|}{axial-only loss} & genet23 \\
$\Delta t$ [ms] & $e_{PV}$ & tw@550 & EDV & $e_{PV}$ & tw@550 & EDV & Fig.~4 $e_{PV}$ \\ \midrule
2    & 2.2e-2 & 14.32 & 136.2 & 7.4e-3 & 5.79 & 136.4 & $\approx$2.4e-3 \\
1    & 2.7e-2 & 15.36 & 135.3 & 1.9e-3 & 5.78 & 136.4 & $\approx$6e-4 \\
0.5  & 2.3e-2 & 15.79 & 133.7 & 3.7e-3 & 5.74 & 136.3 & $\approx$1.8e-4 \\
0.25 & 1.1e-2 & 15.91 & 131.9 & 5.1e-3 & 5.63 & 136.4 & $\approx$4.5e-5 \\
0.125 & ref. & 15.85 & 130.6 & ref. & 5.45 & 136.7 & ref. \\ \bottomrule
\end{tabular}\end{center}
The twist-only loss fails Fig.~4 as badly as the full one: its PV-work error at \SI{1}{ms} is
$45\times$ the paper's, and the EDV drifts by \SI{5.6}{mL}. As $\Delta t\to0$ the numerical dashpot
$2\Delta\mathcal M/\Delta t$ grows, locks the twist and slows filling. \emph{Any} momentum loss
that scales with $1/\Delta t$ is therefore excluded by Fig.~4, whichever components it affects.

\paragraph{Constraints on the missing mechanism.} It must be twist-selective; its damping must
grow with length at least as steeply as $L^3$ (measured $\approx L^{4.4}$); it must not depend on
$\Delta t$ (Fig.~4); and it must be neither a boundary condition (Sec.~2.4, Eq.~44) nor a local
material law. No documented or literature formulation meets all four.

% =========================================================================
\clearpage
\section*{Independent verification of the implementation (no data from genet23)}
Full suite: \url{scripts/ChamberCylinder_verification.py}.\\
Fast subset (pytest): \url{tests/test_chamber_cylinder_verification.py}.\\
The kinematics $r=R+\rho(R)$,
$\theta=\Theta+\beta Z+\varphi(R)$, $z=(1+\varepsilon)Z+\eta(R)$ are the classical
inflation--extension--torsion family of a thick-walled tube. A rotation by $\pi$ about a radial
axis maps the body, the helical fibres (Eq.~22) and the loads onto themselves and $(\varphi,\eta)$
onto $(-\varphi,-\eta)$, so by uniqueness the transmural shears vanish, $\varphi'=\eta'=0$, at all
times. With exact incompressibility, $r^2=r_i^2+(R^2-R_i^2)/\lambda$, the continuum problem then
reduces to three generalized coordinates $q=(r_i,\lambda,\beta)$ plus the active internal
variables. Only the model \emph{definition} (Eqs.~8--36) is shared with the block; kinematics,
quadrature (Gauss--Legendre in $R$), mass, forces and time integration are written independently.
\begin{center}\small\begin{tabular}{P{3.2cm}P{6.0cm}P{6.0cm}}
\toprule
test & reference & agreement \\ \midrule
T1 symmetry & $\varphi'=\eta'=0$; $\alpha\to-\alpha$ flips $\beta$ only & shears $<10^{-14}$; mirror bit-identical \\
T2 static inflation, extension, torsion (passive, $P=0.5$--2\,kPa) & semi-analytical: virtual work
$\partial U/\partial q=P\,\partial V/\partial q$ with exact incompressibility & second-order
convergence; volume $1$--$4\times10^{-4}$, twist $0.002$--$0.005^\circ$ at 24 elements \\
T3 free vibration (isotropic tube, $P=0$) & torsion: $M=\varrho_0J_pL^3/12$, $K=2C_1LJ_p$,
$D=\gamma LJ_p/2+\mu\!\int(\sin\alpha\cos\alpha R)^2\mathrm d\Omega$; axial and breathing modes:
independent modal analysis & frequency $10^{-4}$, decay $2\times10^{-3}$ (torsion, two damping
splits); both axial modes $10^{-3}$; mm/kPa units identical to $5\times10^{-11}$ \\
T4 full beat (Fig.~5 set-up) & independent closed loop: Lagrangian $q$ dynamics, BCS ODEs, valves
and Windkessel (Eqs.~35--36), Radau & ne\,=\,12: volume 0.06\,mL, pressure 0.05\,kPa, torsion
$0.014^\circ$ (max over the beat); tw@550 $5.80^\circ$ vs.\ $5.80^\circ$ \\ \bottomrule
\end{tabular}\end{center}
\begin{center}\includegraphics[width=\linewidth]{ChamberCylinder_verification.png}\end{center}
\begin{itemize}
\item The block solves the written continuum model correctly, including the twist dynamics
  that are in question (inertia with momentum-free rigid rotation, shear stiffness, both
  dampers). An implementation sharing no code with it gives the same fast untwist,
  $5.80^\circ$ at 550\,ms against the paper's $13.3^\circ$.
\item T3 also isolates a damping term that is easy to overlook: with no activation the active
  branch (Eq.~32, $k_s=10^8$) acts as a viscous dashpot $\mu$ on the fibre strain rate. At the
  Table~1 values ($\gamma=\mu=70$) it adds 31\% to the twist damping of Eq.~28. Both
  implementations include it.
\end{itemize}

% =========================================================================
\section*{Discrepancies}
The reproduction matches the paper on pressure, volume and the PV loop across every case
(Figs.~5--9). The following remain, flagged for the authors:
\begin{itemize}
  \item \textbf{Torsion untwist.} With the model now identical to the paper's continuous
  equations (audit \S A), the twist recovers far too early (Fig.~5 tw@550 $5.8^\circ$ vs.\
  $13.3^\circ$), and the residual twist at 600\,ms is independent of the aspect ratio, whereas
  the paper's grows with it (Fig.~6). This is converged in time (to $0.01^\circ$) and in space
  (to $0.05^\circ$), and it survives the full genet23 damping model (tissue viscosity, Eq.~28;
  fibre dissipation, Eq.~32; actin--myosin turnover, Eq.~33) and the full inertia (Eqs.~8, 10,
  18, A8, with free rigid modes). Of the paper's internal inconsistencies, Eq.~59, the missing
  metric in A7 and the $\alpha$-free internal-variable update of Ref.~6 break pressure and volume; the $Z$-independent velocity memory of A7/A9 has the
  right signature but is too strong at the paper's \SI{1}{ms} and, even restricted to the twist,
  contradicts the paper's own Fig.~4. None of the 30 alternative active, passive and viscous laws from the group's and
  colleagues' papers closes the gap either (literature screen). Probing with rotational dashpots
  shows that the figures need an extra twist-selective damping of $\approx3.3\times$ the twist
  damping of Eq.~28, growing with cylinder length as $\approx L^{4.4}$. That scaling excludes every
  local material law, and the paper rules out boundary damping (Sec.~2.4, Eq.~44), so the
  damping has no documented source. An independent re-implementation of the written model gives
  the same torsion as ours (to $0.014^\circ$), so the gap most likely lies between the written
  model and the CardiacLab code that produced the figures. Questions for the authors: how is $\boldsymbol v$
  stored across $Z$; are the rigid-body pins released in dynamics; and which $\Delta t$ and
  start-up produced Fig.~5.
  \item \textbf{Unspecified inputs.} $P_\text{vs}$ is not given in the paper and is fitted to the
  Fig.~5 ESV (\SI{350}{Pa}); $\nu(t)$, $n_0(e_c)$ and $P_\text{at}(t)$ are taken from the authors'
  code (Table~1 refers to Ref.~6).
  \item \textbf{Valve capacitance $C_\text{valve}$.} The valve law reproduces Eq.~35 exactly, but
  the tabulated $C_\text{valve}{=}\SI{9e-9}{m^3/Pa}$ (Eq.~36) is a very large ventricular-node
  compliance ($\approx$\SI{1.2}{mL/mmHg}, $43\times C_\text{ar}$) that stores $\sim$\SI{90}{mL} over
  the pressure swing: applied faithfully it collapses the isovolumic phase (ESV 39 vs.\ 74, peak
  $P$ 8.5 vs.\ 12.8), which is inconsistent with the paper's own Fig.~5. Only
  $C_\text{valve}\le$\,9e-11 is negligible, so we run $C_\text{valve}{\approx}0$; the tabulated
  value appears to be a paper inconsistency.
\end{itemize}

\end{document}
